% Options for packages loaded elsewhere
\PassOptionsToPackage{unicode}{hyperref}
\PassOptionsToPackage{hyphens}{url}
\documentclass[
  11pt,
  a4paper,
]{article}
\usepackage{xcolor}
\usepackage[margin=18mm]{geometry}
\usepackage{amsmath,amssymb}
\setcounter{secnumdepth}{-\maxdimen} % remove section numbering
\usepackage{iftex}
\ifPDFTeX
  \usepackage[T1]{fontenc}
  \usepackage[utf8]{inputenc}
  \usepackage{textcomp} % provide euro and other symbols
\else % if luatex or xetex
  \usepackage{unicode-math} % this also loads fontspec
  \defaultfontfeatures{Scale=MatchLowercase}
  \defaultfontfeatures[\rmfamily]{Ligatures=TeX,Scale=1}
\fi
\ifPDFTeX\else
  % xetex/luatex font selection
    \setmainfont[]{Noto Serif}
    \setsansfont[]{Noto Sans}
    \setmonofont[]{Noto Sans Mono}
  \setmathfont[]{Noto Sans Math}
\fi
% Use upquote if available, for straight quotes in verbatim environments
\IfFileExists{upquote.sty}{\usepackage{upquote}}{}
\IfFileExists{microtype.sty}{% use microtype if available
  \usepackage[]{microtype}
  \UseMicrotypeSet[protrusion]{basicmath} % disable protrusion for tt fonts
}{}
\makeatletter
\@ifundefined{KOMAClassName}{% if non-KOMA class
  \IfFileExists{parskip.sty}{%
    \usepackage{parskip}
  }{% else
    \setlength{\parindent}{0pt}
    \setlength{\parskip}{6pt plus 2pt minus 1pt}}
}{% if KOMA class
  \KOMAoptions{parskip=half}}
\makeatother
\ifLuaTeX
\usepackage[bidi=basic]{babel}
\else
\usepackage[bidi=default]{babel}
\fi
\babelprovide[main,import]{russian}
\ifPDFTeX
\else
\babelfont{rm}[]{Noto Serif}
\fi
% get rid of language-specific shorthands (see #6817):
\let\LanguageShortHands\languageshorthands
\def\languageshorthands#1{}
\setlength{\emergencystretch}{3em} % prevent overfull lines
\providecommand{\tightlist}{%
  \setlength{\itemsep}{0pt}\setlength{\parskip}{0pt}}
\usepackage{bookmark}
\IfFileExists{xurl.sty}{\usepackage{xurl}}{} % add URL line breaks if available
\urlstyle{same}
\hypersetup{
  pdftitle={Билеты по дискретным случайным процессам},
  pdflang={ru-RU},
  hidelinks,
  pdfcreator={LaTeX via pandoc}}

\title{Билеты по дискретным случайным процессам}
\author{}
\date{}

\begin{document}
\maketitle

\section{Билет 17. Стационарные и эргодические случайные
последовательности в широком и узком
смыслах}\label{ux431ux438ux43bux435ux442-17.-ux441ux442ux430ux446ux438ux43eux43dux430ux440ux43dux44bux435-ux438-ux44dux440ux433ux43eux434ux438ux447ux435ux441ux43aux438ux435-ux441ux43bux443ux447ux430ux439ux43dux44bux435-ux43fux43eux441ux43bux435ux434ux43eux432ux430ux442ux435ux43bux44cux43dux43eux441ux442ux438-ux432-ux448ux438ux440ux43eux43aux43eux43c-ux438-ux443ux437ux43aux43eux43c-ux441ux43cux44bux441ux43bux430ux445}

Источники: Ширяев 2, гл. V-VI; лекции ДСП.

\subsection{1. Короткий
ответ}\label{ux43aux43eux440ux43eux442ux43aux438ux439-ux43eux442ux432ux435ux442}

Случайная последовательность

\[
\{X_n:n\in\mathbb Z\}
\]

называется стационарной в узком смысле, если ее конечномерные
распределения не меняются при сдвиге времени:

\[
(X_{n_1+h},\ldots,X_{n_k+h})
\stackrel d=
(X_{n_1},\ldots,X_{n_k})
\]

для всех \(k\), \(n_1,\ldots,n_k\) и целых \(h\).

Стационарность в широком смысле требует только существования вторых
моментов, постоянного среднего и ковариации, зависящей только от лага:

\[
\mathbb E X_n=m,
\qquad
\operatorname{Cov}(X_n,X_{n+k})=R(k).
\]

Эргодичность в узком смысле означает, что сдвиг-инвариантные события
имеют вероятность только \(0\) или \(1\). Интуитивно: одна достаточно
длинная траектория отражает вероятностные свойства процесса.

Эргодичность в широком смысле обычно понимают как сходимость временного
среднего к математическому ожиданию в среднем квадратичном:

\[
\frac1N\sum_{n=0}^{N-1}X_n
\xrightarrow[L^2]{N\to\infty}
m.
\]

\subsection{2. Полный
ответ}\label{ux43fux43eux43bux43dux44bux439-ux43eux442ux432ux435ux442}

\subsubsection{17.1. Введение и основная
идея}\label{ux432ux432ux435ux434ux435ux43dux438ux435-ux438-ux43eux441ux43dux43eux432ux43dux430ux44f-ux438ux434ux435ux44f}

Этот билет продолжает Билет 16, где дискретный процесс рассматривался
как случайная последовательность и как случайный элемент пространства
последовательностей. Теперь нас интересуют свойства последовательности,
которые описывают поведение при сдвиге времени.

Пусть

\[
X=(X_n)_{n\in\mathbb Z}
\]

или

\[
X=(X_n)_{n\ge0}
\]

\begin{itemize}
\tightlist
\item
  случайная последовательность. Для стационарности удобнее работать с
  двусторонним временем \(n\in\mathbb Z\), потому что сдвиг можно
  применять в обе стороны. Для экзамена можно формулировать и для
  \(n\ge0\), но тогда сдвиги берут только такие, чтобы индексы
  оставались неотрицательными.
\end{itemize}

Главная идея стационарности: статистические свойства процесса не
меняются при переносе начала отсчета времени. Если сегодня процесс
выглядит так же, как завтра, в смысле распределений, то он стационарен.

\subsubsection{17.2. Стационарность в узком
смысле}\label{ux441ux442ux430ux446ux438ux43eux43dux430ux440ux43dux43eux441ux442ux44c-ux432-ux443ux437ux43aux43eux43c-ux441ux43cux44bux441ux43bux435}

Сначала определим стационарность в узком смысле. Последовательность
называется стационарной в узком смысле, или строго стационарной, если
для любых

\[
k\ge1,\qquad n_1,\ldots,n_k\in\mathbb Z,\qquad h\in\mathbb Z
\]

выполнено равенство распределений:

\[
(X_{n_1+h},\ldots,X_{n_k+h})
\stackrel d=
(X_{n_1},\ldots,X_{n_k}).
\]

Это означает, что все конечномерные распределения инвариантны
относительно сдвига времени. Например,

\[
X_0\stackrel d=X_{10},
\]

но этого мало: также должно быть

\[
(X_0,X_1,X_5)\stackrel d=(X_{10},X_{11},X_{15})
\]

и аналогично для любого конечного набора моментов.

На языке пространства последовательностей из Билета 16 стационарность в
узком смысле можно выразить через сдвиг. Пусть

\[
S:E^{\mathbb Z}\to E^{\mathbb Z}
\]

задан формулой

\[
(Sx)_n=x_{n+1}.
\]

Если \(P_X\) - закон последовательности на \(E^{\mathbb Z}\), то узкая
стационарность означает

\[
P_X\circ S^{-1}=P_X.
\]

То есть закон процесса как мера на пространстве траекторий инвариантен
относительно сдвига.

\subsubsection{17.3. Стационарность в широком
смысле}\label{ux441ux442ux430ux446ux438ux43eux43dux430ux440ux43dux43eux441ux442ux44c-ux432-ux448ux438ux440ux43eux43aux43eux43c-ux441ux43cux44bux441ux43bux435}

Теперь стационарность в широком смысле. Она относится только к
характеристикам второго порядка, поэтому требует

\[
\mathbb E X_n^2<\infty
\]

для всех \(n\). Последовательность стационарна в широком смысле, если:

\begin{enumerate}
\def\labelenumi{\arabic{enumi}.}
\tightlist
\item
  среднее не зависит от времени:
\end{enumerate}

\[
\mathbb E X_n=m;
\]

\begin{enumerate}
\def\labelenumi{\arabic{enumi}.}
\setcounter{enumi}{1}
\tightlist
\item
  ковариация зависит только от разности индексов:
\end{enumerate}

\[
\operatorname{Cov}(X_n,X_m)=R(m-n).
\]

Часто записывают

\[
R(k)=\operatorname{Cov}(X_n,X_{n+k}).
\]

Если последовательность центрирована, то есть \(m=0\), то

\[
R(k)=\mathbb E[X_nX_{n+k}].
\]

Иногда в курсах случайных процессов корреляционной функцией называют
именно

\[
\mathbb E[X_nX_{n+k}],
\]

а ковариационной - центрированную величину. Важно проговорить, какая
версия используется. В этих билетах удобнее держать в голове:

\[
\operatorname{Cov}(X_n,X_{n+k})
=
\mathbb E[(X_n-m)(X_{n+k}-m)].
\]

\subsubsection{17.4. Связь между типами
стационарности}\label{ux441ux432ux44fux437ux44c-ux43cux435ux436ux434ux443-ux442ux438ux43fux430ux43cux438-ux441ux442ux430ux446ux438ux43eux43dux430ux440ux43dux43eux441ux442ux438}

Узкая стационарность при существовании вторых моментов влечет широкую
стационарность. Действительно, из равенства одномерных распределений
следует одинаковое среднее:

\[
\mathbb E X_n=m.
\]

Из равенства двумерных распределений

\[
(X_n,X_{n+k})\stackrel d=(X_0,X_k)
\]

следует

\[
\operatorname{Cov}(X_n,X_{n+k})
=
\operatorname{Cov}(X_0,X_k).
\]

Обратное в общем случае неверно: широкая стационарность контролирует
только первые два момента, но не все распределение.

Есть важное исключение: для гауссовских последовательностей широкая
стационарность влечет узкую. Причина в том, что конечномерные
распределения гауссовского процесса полностью задаются средним вектором
и ковариационной матрицей. Если среднее постоянно, а ковариация зависит
только от лага, то все конечномерные гауссовские распределения
инвариантны при сдвиге. Это связь с Билетом 18.

\subsubsection{17.5. Понятие эргодичности и инвариантные
события}\label{ux43fux43eux43dux44fux442ux438ux435-ux44dux440ux433ux43eux434ux438ux447ux43dux43eux441ux442ux438-ux438-ux438ux43dux432ux430ux440ux438ux430ux43dux442ux43dux44bux435-ux441ux43eux431ux44bux442ux438ux44f}

Теперь эргодичность. Стационарность говорит, что закон не меняется при
сдвиге. Эргодичность говорит сильнее: в стационарном процессе нет
скрытой случайной характеристики, которая сохраняется при сдвиге и
разбивает траектории на разные режимы.

На пространстве траекторий событие \(A\) называется сдвиг-инвариантным,
если

\[
S^{-1}A=A
\]

с точностью до множеств вероятности ноль. Интуитивно такое событие не
меняется, если сдвинуть всю последовательность на один шаг.

Стационарная последовательность называется эргодической в узком смысле,
если любое сдвиг-инвариантное событие имеет вероятность \(0\) или \(1\):

\[
A=S^{-1}A
\quad\Longrightarrow\quad
P_X(A)\in\{0,1\}.
\]

Это определение часто называют метрической эргодичностью. Его смысл:
сдвиг не оставляет нетривиальных случайных ``меток'' траектории.

\subsubsection{17.6. Практическая интерпретация и временные
средние}\label{ux43fux440ux430ux43aux442ux438ux447ux435ux441ux43aux430ux44f-ux438ux43dux442ux435ux440ux43fux440ux435ux442ux430ux446ux438ux44f-ux438-ux432ux440ux435ux43cux435ux43dux43dux44bux435-ux441ux440ux435ux434ux43dux438ux435}

Практическая интерпретация эргодичности связана с временными средними.
Если процесс стационарен и эргодичен в узком смысле, то по эргодической
теореме для любой измеримой функции \(f\) с
\(\mathbb E|f(X_0)|<\infty\):

\[
\frac1N\sum_{n=0}^{N-1} f(X_n)
\to
\mathbb E f(X_0)
\]

почти наверное. В частности, для \(f(x)=x\), если
\(\mathbb E|X_0|<\infty\):

\[
\frac1N\sum_{n=0}^{N-1}X_n
\to
\mathbb E X_0
\]

почти наверное.

Это и есть главный прикладной смысл: среднее по одной длинной реализации
можно использовать вместо среднего по ансамблю.

\subsubsection{17.7. Эргодичность в широком
смысле}\label{ux44dux440ux433ux43eux434ux438ux447ux43dux43eux441ux442ux44c-ux432-ux448ux438ux440ux43eux43aux43eux43c-ux441ux43cux44bux441ux43bux435}

Эргодичность в широком смысле - более слабое понятие второго порядка.
Для стационарной в широком смысле последовательности с конечными вторыми
моментами говорят, что она эргодична по среднему, или эргодична в
широком смысле относительно среднего, если

\[
\overline X_N
=
\frac1N\sum_{n=0}^{N-1}X_n
\]

сходится к \(m=\mathbb EX_0\) в среднем квадратичном:

\[
\mathbb E|\overline X_N-m|^2\to0.
\]

Эквивалентно,

\[
\operatorname{Var}(\overline X_N)\to0.
\]

Для стационарной в широком смысле последовательности

\[
\operatorname{Var}(\overline X_N)
=
\frac1{N^2}
\sum_{i=0}^{N-1}\sum_{j=0}^{N-1}R(j-i),
\]

где

\[
R(k)=\operatorname{Cov}(X_0,X_k).
\]

Поэтому условия затухания корреляций дают эргодичность среднего. Это
подробно развивается в Билете 24.

\subsubsection{17.8. Сравнение понятий и
примеры}\label{ux441ux440ux430ux432ux43dux435ux43dux438ux435-ux43fux43eux43dux44fux442ux438ux439-ux438-ux43fux440ux438ux43cux435ux440ux44b}

Важно различать:

\begin{enumerate}
\def\labelenumi{\arabic{enumi}.}
\tightlist
\item
  стационарность - инвариантность распределений или моментов при сдвиге;
\item
  эргодичность - возможность восстанавливать вероятностные
  характеристики по одной длинной реализации;
\item
  независимость - совсем другое свойство, гораздо сильнее во многих
  примерах.
\end{enumerate}

Например, iid-последовательность стационарна и эргодична, но
стационарность сама по себе не означает независимость.
Последовательность \(X_n=Y\), где \(Y\) - одна случайная величина,
стационарна, но не эргодична: траектория хранит неизменный случайный
уровень \(Y\).

\subsection{3. Основные
определения}\label{ux43eux441ux43dux43eux432ux43dux44bux435-ux43eux43fux440ux435ux434ux435ux43bux435ux43dux438ux44f}

\subsubsection{Сдвиг на пространстве
последовательностей}\label{ux441ux434ux432ux438ux433-ux43dux430-ux43fux440ux43eux441ux442ux440ux430ux43dux441ux442ux432ux435-ux43fux43eux441ux43bux435ux434ux43eux432ux430ux442ux435ux43bux44cux43dux43eux441ux442ux435ux439}

Для двусторонних последовательностей \(x=(x_n)_{n\in\mathbb Z}\):

\[
(Sx)_n=x_{n+1}.
\]

\subsubsection{Стационарность в узком
смысле}\label{ux441ux442ux430ux446ux438ux43eux43dux430ux440ux43dux43eux441ux442ux44c-ux432-ux443ux437ux43aux43eux43c-ux441ux43cux44bux441ux43bux435-1}

Последовательность \((X_n)\) стационарна в узком смысле, если для любых
\(k\), \(n_1,\ldots,n_k\) и \(h\):

\[
(X_{n_1+h},\ldots,X_{n_k+h})
\stackrel d=
(X_{n_1},\ldots,X_{n_k}).
\]

\subsubsection{Стационарность в широком
смысле}\label{ux441ux442ux430ux446ux438ux43eux43dux430ux440ux43dux43eux441ux442ux44c-ux432-ux448ux438ux440ux43eux43aux43eux43c-ux441ux43cux44bux441ux43bux435-1}

Последовательность с конечными вторыми моментами стационарна в широком
смысле, если

\[
\mathbb E X_n=m
\]

не зависит от \(n\), а

\[
\operatorname{Cov}(X_n,X_m)=R(m-n)
\]

зависит только от лага.

\subsubsection{Корреляционная или ковариационная функция стационарной
последовательности}\label{ux43aux43eux440ux440ux435ux43bux44fux446ux438ux43eux43dux43dux430ux44f-ux438ux43bux438-ux43aux43eux432ux430ux440ux438ux430ux446ux438ux43eux43dux43dux430ux44f-ux444ux443ux43dux43aux446ux438ux44f-ux441ux442ux430ux446ux438ux43eux43dux430ux440ux43dux43eux439-ux43fux43eux441ux43bux435ux434ux43eux432ux430ux442ux435ux43bux44cux43dux43eux441ux442ux438}

Для центрированной последовательности:

\[
R(k)=\mathbb E[X_0X_k].
\]

В общем случае:

\[
R(k)=\operatorname{Cov}(X_0,X_k).
\]

\subsubsection{Сдвиг-инвариантное
событие}\label{ux441ux434ux432ux438ux433-ux438ux43dux432ux430ux440ux438ux430ux43dux442ux43dux43eux435-ux441ux43eux431ux44bux442ux438ux435}

Событие \(A\) на пространстве траекторий такое, что

\[
S^{-1}A=A
\]

с точностью до множеств вероятности ноль.

\subsubsection{Эргодичность в узком
смысле}\label{ux44dux440ux433ux43eux434ux438ux447ux43dux43eux441ux442ux44c-ux432-ux443ux437ux43aux43eux43c-ux441ux43cux44bux441ux43bux435}

Стационарная последовательность эргодична, если каждое
сдвиг-инвариантное событие имеет вероятность \(0\) или \(1\):

\[
S^{-1}A=A
\quad\Rightarrow\quad
P(A)\in\{0,1\}.
\]

\subsubsection{Эргодичность в широком смысле по
среднему}\label{ux44dux440ux433ux43eux434ux438ux447ux43dux43eux441ux442ux44c-ux432-ux448ux438ux440ux43eux43aux43eux43c-ux441ux43cux44bux441ux43bux435-ux43fux43e-ux441ux440ux435ux434ux43dux435ux43cux443}

Стационарная в широком смысле последовательность эргодична по среднему,
если

\[
\frac1N\sum_{n=0}^{N-1}X_n
\xrightarrow[L^2]{}
\mathbb E X_0.
\]

\subsection{4. Основные теоремы и
свойства}\label{ux43eux441ux43dux43eux432ux43dux44bux435-ux442ux435ux43eux440ux435ux43cux44b-ux438-ux441ux432ux43eux439ux441ux442ux432ux430}

\begin{itemize}
\tightlist
\item
  Узкая стационарность означает инвариантность закона процесса
  относительно сдвига:
\end{itemize}

\[
P_X\circ S^{-1}=P_X.
\]

\begin{itemize}
\tightlist
\item
  Узкая стационарность и конечные вторые моменты влекут широкую
  стационарность.
\item
  Широкая стационарность в общем случае не влечет узкую.
\item
  Для гауссовской последовательности широкая стационарность влечет
  узкую.
\item
  Эргодичность в узком смысле равносильна отсутствию нетривиальных
  сдвиг-инвариантных событий.
\item
  Для стационарной эргодической последовательности временные средние
  интегрируемых функций совпадают с математическими ожиданиями:
\end{itemize}

\[
\frac1N\sum_{n=0}^{N-1}f(X_n)\to\mathbb E f(X_0)
\]

почти наверное.

\begin{itemize}
\tightlist
\item
  Для стационарной в широком смысле последовательности эргодичность
  среднего связана с условием:
\end{itemize}

\[
\operatorname{Var}\left(\frac1N\sum_{n=0}^{N-1}X_n\right)\to0.
\]

\begin{itemize}
\tightlist
\item
  Достаточно, например, чтобы
\end{itemize}

\[
\frac1N\sum_{k=1}^{N}|R(k)|\to0.
\]

В частности, если \(R(k)\to0\), то временное среднее сходится к
математическому ожиданию в \(L^2\). Общий критерий второго порядка:

\[
\operatorname{Var}(\overline X_N)=
\frac1{N^2}\sum_{i,j=0}^{N-1}R(j-i)\to0.
\]

\subsection{5. Примеры}\label{ux43fux440ux438ux43cux435ux440ux44b}

\subsubsection{iid-последовательность}\label{iid-ux43fux43eux441ux43bux435ux434ux43eux432ux430ux442ux435ux43bux44cux43dux43eux441ux442ux44c}

Пусть

\[
X_0,X_1,\ldots
\]

независимы и одинаково распределены.

Тогда последовательность стационарна в узком смысле: любой конечный
вектор имеет произведение одних и тех же одномерных законов, и сдвиг
времени не меняет это произведение.

Если \(\mathbb E X_0^2<\infty\), то она стационарна и в широком смысле:

\[
\mathbb E X_n=m,
\]

и

\[
\operatorname{Cov}(X_n,X_{n+k})=
\begin{cases}
\operatorname{Var}X_0, & k=0,\\
0, & k\ne0.
\end{cases}
\]

Такая последовательность эргодична: временные средние сходятся к
математическому ожиданию по закону больших чисел.

\subsubsection{Случайная
константа}\label{ux441ux43bux443ux447ux430ux439ux43dux430ux44f-ux43aux43eux43dux441ux442ux430ux43dux442ux430}

Пусть

\[
X_n=Y
\]

для всех \(n\), где \(Y\) - неконстантная случайная величина.

Эта последовательность стационарна в узком смысле: сдвиг времени вообще
ничего не меняет. Если \(\mathbb E Y^2<\infty\), она стационарна и в
широком смысле:

\[
R(k)=\operatorname{Var}Y
\]

для всех \(k\).

Но она не эргодична. Временное среднее равно

\[
\frac1N\sum_{n=0}^{N-1}X_n=Y,
\]

а не константе \(\mathbb EY\). Одна траектория показывает только свой
случайный уровень.

\subsubsection{Широкая стационарность без
узкой}\label{ux448ux438ux440ux43eux43aux430ux44f-ux441ux442ux430ux446ux438ux43eux43dux430ux440ux43dux43eux441ux442ux44c-ux431ux435ux437-ux443ux437ux43aux43eux439}

Пусть величины \(X_n\) независимы, имеют нулевое среднее и единичную
дисперсию, но распределения чередуются:

\[
X_{2k}\sim N(0,1),
\qquad
X_{2k+1}\sim
\begin{cases}
1, & 1/2,\\
-1, & 1/2.
\end{cases}
\]

Тогда

\[
\mathbb E X_n=0,
\qquad
\operatorname{Var}X_n=1,
\]

а при \(n\ne m\):

\[
\operatorname{Cov}(X_n,X_m)=0
\]

из-за независимости. Значит последовательность стационарна в широком
смысле:

\[
R(0)=1,\qquad R(k)=0,\ k\ne0.
\]

Но она не стационарна в узком смысле, потому что одномерное
распределение зависит от четности \(n\).

\subsubsection{Гауссовская стационарная
последовательность}\label{ux433ux430ux443ux441ux441ux43eux432ux441ux43aux430ux44f-ux441ux442ux430ux446ux438ux43eux43dux430ux440ux43dux430ux44f-ux43fux43eux441ux43bux435ux434ux43eux432ux430ux442ux435ux43bux44cux43dux43eux441ux442ux44c}

Если последовательность гауссовская, то ее конечномерные распределения
полностью задаются средними и ковариациями. Поэтому для гауссовской
последовательности условия

\[
\mathbb E X_n=m,
\qquad
\operatorname{Cov}(X_n,X_{n+k})=R(k)
\]

уже дают узкую стационарность.

\subsubsection{Гармоника со случайной
фазой}\label{ux433ux430ux440ux43cux43eux43dux438ux43aux430-ux441ux43e-ux441ux43bux443ux447ux430ux439ux43dux43eux439-ux444ux430ux437ux43eux439}

Пусть

\[
X_n=\cos(\lambda n+\Phi),
\]

где \(\Phi\) равномерна на \([0,2\pi]\). Тогда сдвиг \(n\mapsto n+h\)
заменяет фазу \(\Phi\) на \(\Phi+\lambda h\), а равномерное
распределение фазы не меняется. Поэтому последовательность стационарна в
узком смысле.

Этот пример полезен тем, что зависимость во времени сильная, но
распределения при сдвиге не меняются.

\subsection{6. Схемы доказательств /
обоснований}\label{ux441ux445ux435ux43cux44b-ux434ux43eux43aux430ux437ux430ux442ux435ux43bux44cux441ux442ux432-ux43eux431ux43eux441ux43dux43eux432ux430ux43dux438ux439}

\subsubsection{Почему узкая стационарность влечет
широкую}\label{ux43fux43eux447ux435ux43cux443-ux443ux437ux43aux430ux44f-ux441ux442ux430ux446ux438ux43eux43dux430ux440ux43dux43eux441ux442ux44c-ux432ux43bux435ux447ux435ux442-ux448ux438ux440ux43eux43aux443ux44e}

Пусть последовательность стационарна в узком смысле и

\[
\mathbb E X_n^2<\infty.
\]

Из

\[
X_n\stackrel d=X_0
\]

следует

\[
\mathbb E X_n=\mathbb E X_0=m.
\]

Из

\[
(X_n,X_{n+k})\stackrel d=(X_0,X_k)
\]

следует равенство вторых смешанных моментов:

\[
\mathbb E[X_nX_{n+k}]=\mathbb E[X_0X_k].
\]

Значит

\[
\operatorname{Cov}(X_n,X_{n+k})
=
\operatorname{Cov}(X_0,X_k),
\]

то есть ковариация зависит только от лага.

\subsubsection{Почему широкая стационарность не влечет
узкую}\label{ux43fux43eux447ux435ux43cux443-ux448ux438ux440ux43eux43aux430ux44f-ux441ux442ux430ux446ux438ux43eux43dux430ux440ux43dux43eux441ux442ux44c-ux43dux435-ux432ux43bux435ux447ux435ux442-ux443ux437ux43aux443ux44e}

Широкая стационарность фиксирует только

\[
\mathbb E X_n
\]

и

\[
\operatorname{Cov}(X_n,X_m).
\]

Она не видит третьи, четвертые и вообще все более высокие
распределительные характеристики. Поэтому можно сделать
последовательность с одинаковыми средними и ковариациями, но с разными
одномерными распределениями в разные моменты времени. Пример с
чередованием нормального и радемахеровского распределения показывает это
явно.

\subsubsection{Почему для гауссовских широкая стационарность дает
узкую}\label{ux43fux43eux447ux435ux43cux443-ux434ux43bux44f-ux433ux430ux443ux441ux441ux43eux432ux441ux43aux438ux445-ux448ux438ux440ux43eux43aux430ux44f-ux441ux442ux430ux446ux438ux43eux43dux430ux440ux43dux43eux441ux442ux44c-ux434ux430ux435ux442-ux443ux437ux43aux443ux44e}

Гауссовский вектор полностью определяется средним вектором и
ковариационной матрицей. Для любого набора моментов

\[
n_1,\ldots,n_k
\]

средний вектор равен

\[
(m,\ldots,m),
\]

а ковариационная матрица имеет элементы

\[
\operatorname{Cov}(X_{n_i},X_{n_j})
=
R(n_j-n_i).
\]

После сдвига всех индексов на \(h\) разности не меняются:

\[
(n_j+h)-(n_i+h)=n_j-n_i.
\]

Значит средний вектор и ковариационная матрица не меняются,
следовательно не меняется и гауссовское конечномерное распределение.

\subsubsection{Почему случайная константа не
эргодична}\label{ux43fux43eux447ux435ux43cux443-ux441ux43bux443ux447ux430ux439ux43dux430ux44f-ux43aux43eux43dux441ux442ux430ux43dux442ux430-ux43dux435-ux44dux440ux433ux43eux434ux438ux447ux43dux430}

Если \(X_n=Y\), то событие

\[
A=\{Y>0\}
\]

не меняется при сдвиге траектории. Если

\[
0<P\{Y>0\}<1,
\]

то это нетривиальное сдвиг-инвариантное событие. Значит процесс не
эргодичен.

\subsection{7. Что написать на
доске}\label{ux447ux442ux43e-ux43dux430ux43fux438ux441ux430ux442ux44c-ux43dux430-ux434ux43eux441ux43aux435}

\begin{enumerate}
\def\labelenumi{\arabic{enumi}.}
\tightlist
\item
  Узкая стационарность:
\end{enumerate}

\[
(X_{n_1+h},\ldots,X_{n_k+h})
\stackrel d=
(X_{n_1},\ldots,X_{n_k}).
\]

\begin{enumerate}
\def\labelenumi{\arabic{enumi}.}
\setcounter{enumi}{1}
\tightlist
\item
  Сдвиг:
\end{enumerate}

\[
(Sx)_n=x_{n+1},
\qquad
P_X\circ S^{-1}=P_X.
\]

\begin{enumerate}
\def\labelenumi{\arabic{enumi}.}
\setcounter{enumi}{2}
\tightlist
\item
  Широкая стационарность:
\end{enumerate}

\[
\mathbb E X_n=m,
\qquad
\operatorname{Cov}(X_n,X_{n+k})=R(k).
\]

\begin{enumerate}
\def\labelenumi{\arabic{enumi}.}
\setcounter{enumi}{3}
\tightlist
\item
  Узкая плюс вторые моменты:
\end{enumerate}

\[
\text{узкая стационарность}+\mathbb E X_n^2<\infty
\Longrightarrow
\text{широкая стационарность}.
\]

\begin{enumerate}
\def\labelenumi{\arabic{enumi}.}
\setcounter{enumi}{4}
\tightlist
\item
  Эргодичность:
\end{enumerate}

\[
S^{-1}A=A
\Longrightarrow
P(A)\in\{0,1\}.
\]

\begin{enumerate}
\def\labelenumi{\arabic{enumi}.}
\setcounter{enumi}{5}
\tightlist
\item
  Эргодичность среднего:
\end{enumerate}

\[
\frac1N\sum_{n=0}^{N-1}X_n
\to
\mathbb E X_0.
\]

Для широкого смысла обычно:

\[
\frac1N\sum_{n=0}^{N-1}X_n
\xrightarrow{L^2}
\mathbb E X_0.
\]

\subsection{8. Типичные дополнительные
вопросы}\label{ux442ux438ux43fux438ux447ux43dux44bux435-ux434ux43eux43fux43eux43bux43dux438ux442ux435ux43bux44cux43dux44bux435-ux432ux43eux43fux440ux43eux441ux44b}

\textbf{Что является главным объектом билета?}

Случайная последовательность \((X_n)\) и ее поведение при сдвиге
времени. Нужно различать инвариантность всех конечномерных
распределений, инвариантность первых двух моментов и эргодичность.

\textbf{В чем разница между узкой и широкой стационарностью?}

Узкая стационарность требует равенства всех конечномерных распределений
после сдвига:

\[
(X_{n_1+h},\ldots,X_{n_k+h})
\stackrel d=
(X_{n_1},\ldots,X_{n_k}).
\]

Широкая стационарность требует только постоянного среднего и ковариации,
зависящей от лага:

\[
\mathbb E X_n=m,
\qquad
\operatorname{Cov}(X_n,X_{n+k})=R(k).
\]

\textbf{Почему узкая стационарность сильнее широкой?}

Потому что она фиксирует все конечномерные распределения. Если есть
вторые моменты, из этих распределений следуют и средние, и ковариации.
Широкая стационарность фиксирует только первые два момента.

\textbf{Когда широкая стационарность влечет узкую?}

Для гауссовских последовательностей. Гауссовские конечномерные
распределения полностью задаются средним вектором и ковариационной
матрицей.

\textbf{Что такое эргодичность в узком смысле?}

Это отсутствие нетривиальных сдвиг-инвариантных событий:

\[
S^{-1}A=A
\Rightarrow
P(A)\in\{0,1\}.
\]

Интуитивно одна длинная траектория достаточно хорошо представляет весь
ансамбль.

\textbf{Что такое эргодичность в широком смысле?}

Обычно в этом курсе это эргодичность среднего: временное среднее
сходится к математическому ожиданию в \(L^2\):

\[
\frac1N\sum_{n=0}^{N-1}X_n
\xrightarrow{L^2}
\mathbb E X_0.
\]

\textbf{Стационарность означает независимость?}

Нет. Стационарность говорит об инвариантности распределений при сдвиге,
но не запрещает зависимость во времени. Например, \(X_n=Y\) для всех
\(n\) стационарна, но координаты полностью зависимы.

\textbf{Эргодичность означает независимость?}

Нет. Независимые одинаково распределенные последовательности обычно
эргодичны, но эргодичность не равна независимости. Эргодичность
запрещает сдвиг-инвариантные случайные режимы, но допускает зависимость.

\textbf{Почему случайная константа стационарна, но не эргодична?}

Потому что сдвиг не меняет траекторию, значит стационарность есть. Но
временное среднее равно случайной величине \(Y\), а не константе
\(\mathbb EY\). Кроме того, событие \(\{Y>0\}\) сдвиг-инвариантно и
может иметь вероятность между \(0\) и \(1\).

\textbf{Как билет связан с Билетом 24?}

В Билете 24 изучается корреляционная функция

\[
R(k)=\operatorname{Cov}(X_0,X_k)
\]

и достаточные условия, при которых временное среднее сходится к
математическому ожиданию в \(L^2\).

\subsection{9. Типичные
ошибки}\label{ux442ux438ux43fux438ux447ux43dux44bux435-ux43eux448ux438ux431ux43aux438}

\begin{itemize}
\tightlist
\item
  Смешивать стационарность и эргодичность.
\item
  Считать, что стационарность означает независимость.
\item
  Считать, что широкая стационарность всегда равносильна узкой.
\item
  Забывать условие существования вторых моментов для широкой
  стационарности.
\item
  Проверять только одномерные распределения и объявлять узкую
  стационарность.
\item
  Называть процесс эргодическим без указания смысла: узкий смысл,
  среднее, \(L^2\).
\item
  Путать ковариационную функцию \(R(k)\) и коэффициент корреляции.
\item
  Забывать, что случайная константа является стационарной, но не
  эргодической.
\end{itemize}

\subsection{10. Связи с другими
билетами}\label{ux441ux432ux44fux437ux438-ux441-ux434ux440ux443ux433ux438ux43cux438-ux431ux438ux43bux435ux442ux430ux43cux438}

\begin{itemize}
\tightlist
\item
  Билет 16: стационарность формулируется для закона на пространстве
  последовательностей.
\item
  Билет 15: узкая стационарность - свойство согласованных конечномерных
  распределений.
\item
  Билет 14: широкая стационарность использует ковариации и геометрию
  второго порядка.
\item
  Билет 18: для гауссовских последовательностей широкая стационарность
  влечет узкую.
\item
  Билет 21: временные средние и закон больших чисел связаны с
  эргодической идеей.
\item
  Билет 24: корреляционная функция и достаточные условия эргодичности
  среднего.
\item
  Билет 25: спектральные характеристики строятся для стационарных
  последовательностей второго порядка.
\end{itemize}

\subsection{11. Минимум на
удовлетворительно}\label{ux43cux438ux43dux438ux43cux443ux43c-ux43dux430-ux443ux434ux43eux432ux43bux435ux442ux432ux43eux440ux438ux442ux435ux43bux44cux43dux43e}

Нужно сказать:

\begin{enumerate}
\def\labelenumi{\arabic{enumi}.}
\tightlist
\item
  Узкая стационарность:
\end{enumerate}

\[
(X_{n_1+h},\ldots,X_{n_k+h})
\stackrel d=
(X_{n_1},\ldots,X_{n_k}).
\]

\begin{enumerate}
\def\labelenumi{\arabic{enumi}.}
\setcounter{enumi}{1}
\tightlist
\item
  Широкая стационарность:
\end{enumerate}

\[
\mathbb E X_n=m,
\qquad
\operatorname{Cov}(X_n,X_{n+k})=R(k).
\]

\begin{enumerate}
\def\labelenumi{\arabic{enumi}.}
\setcounter{enumi}{2}
\item
  Узкая плюс вторые моменты влечет широкую.
\item
  Эргодичность: временные средние отражают вероятностные средние; строго
  - нет нетривиальных сдвиг-инвариантных событий.
\item
  Главная ловушка: стационарность не равна эргодичности и не равна
  независимости.
\end{enumerate}

\subsection{12. Минимум на
хорошо}\label{ux43cux438ux43dux438ux43cux443ux43c-ux43dux430-ux445ux43eux440ux43eux448ux43e}

Помимо минимума, нужно объяснить:

\begin{itemize}
\tightlist
\item
  что такое сдвиг на пространстве последовательностей;
\item
  почему узкая стационарность означает инвариантность закона
  относительно сдвига;
\item
  почему широкая стационарность требует вторых моментов;
\item
  почему широкая не всегда влечет узкую;
\item
  почему для гауссовских последовательностей широкая стационарность
  влечет узкую;
\item
  разобрать пример iid-последовательности и случайной константы.
\end{itemize}

\subsection{13. Что добавить для
отлично}\label{ux447ux442ux43e-ux434ux43eux431ux430ux432ux438ux442ux44c-ux434ux43bux44f-ux43eux442ux43bux438ux447ux43dux43e}

Для сильного ответа стоит добавить:

\begin{itemize}
\tightlist
\item
  формулировку эргодичности через сдвиг-инвариантные события;
\item
  интерпретацию через эргодическую теорему:
\end{itemize}

\[
\frac1N\sum_{n=0}^{N-1}f(X_n)\to\mathbb E f(X_0);
\]

\begin{itemize}
\tightlist
\item
  определение эргодичности в широком смысле как \(L^2\)-сходимости
  среднего;
\item
  формулу для дисперсии среднего:
\end{itemize}

\[
\operatorname{Var}(\overline X_N)
=
\frac1{N^2}\sum_{i,j=0}^{N-1}R(j-i);
\]

\begin{itemize}
\tightlist
\item
  пример широкой стационарности без узкой;
\item
  пример стационарной неэргодической последовательности;
\item
  связь со спектральной теорией и корреляционной функцией.
\end{itemize}

\subsection{14. Устная версия ответа на 5
минут}\label{ux443ux441ux442ux43dux430ux44f-ux432ux435ux440ux441ux438ux44f-ux43eux442ux432ux435ux442ux430-ux43dux430-5-ux43cux438ux43dux443ux442}

\begin{enumerate}
\def\labelenumi{\arabic{enumi}.}
\tightlist
\item
  В этом билете изучаются свойства случайной последовательности при
  сдвиге времени.
\item
  Узкая стационарность означает, что все конечномерные распределения не
  меняются при сдвиге:
\end{enumerate}

\[
(X_{n_1+h},\ldots,X_{n_k+h})
\stackrel d=
(X_{n_1},\ldots,X_{n_k}).
\]

\begin{enumerate}
\def\labelenumi{\arabic{enumi}.}
\setcounter{enumi}{2}
\tightlist
\item
  На пространстве траекторий это означает инвариантность закона
  относительно сдвига:
\end{enumerate}

\[
P_X\circ S^{-1}=P_X.
\]

\begin{enumerate}
\def\labelenumi{\arabic{enumi}.}
\setcounter{enumi}{3}
\tightlist
\item
  Широкая стационарность требует вторых моментов, постоянного среднего и
  ковариации только от лага:
\end{enumerate}

\[
\mathbb E X_n=m,
\qquad
\operatorname{Cov}(X_n,X_{n+k})=R(k).
\]

\begin{enumerate}
\def\labelenumi{\arabic{enumi}.}
\setcounter{enumi}{4}
\tightlist
\item
  Узкая стационарность при наличии вторых моментов влечет широкую.
\item
  Обратное неверно, потому что широкая стационарность контролирует
  только первые два момента.
\item
  Для гауссовских последовательностей обратное верно: среднее и
  ковариация задают все конечномерные распределения.
\item
  Эргодичность в узком смысле означает, что сдвиг-инвариантные события
  имеют вероятность \(0\) или \(1\).
\item
  Практически это означает, что временные средние совпадают с
  вероятностными:
\end{enumerate}

\[
\frac1N\sum_{n=0}^{N-1}f(X_n)\to\mathbb E f(X_0).
\]

\begin{enumerate}
\def\labelenumi{\arabic{enumi}.}
\setcounter{enumi}{9}
\tightlist
\item
  В широком смысле обычно говорят об эргодичности среднего:
\end{enumerate}

\[
\frac1N\sum_{n=0}^{N-1}X_n
\xrightarrow{L^2}
\mathbb E X_0.
\]

\begin{enumerate}
\def\labelenumi{\arabic{enumi}.}
\setcounter{enumi}{10}
\tightlist
\item
  Примеры: iid-последовательность стационарна и эргодична; случайная
  константа стационарна, но не эргодична.
\item
  Завершить можно связью с корреляционной функцией: условия на \(R(k)\)
  дают эргодичность среднего.
\end{enumerate}

\subsection{15. Если осталось 2 минуты перед
экзаменом}\label{ux435ux441ux43bux438-ux43eux441ux442ux430ux43bux43eux441ux44c-2-ux43cux438ux43dux443ux442ux44b-ux43fux435ux440ux435ux434-ux44dux43aux437ux430ux43cux435ux43dux43eux43c}

Запомнить три различия.

Узкая стационарность:

\[
(X_{n_1+h},\ldots,X_{n_k+h})
\stackrel d=
(X_{n_1},\ldots,X_{n_k}).
\]

Широкая стационарность:

\[
\mathbb E X_n=m,
\qquad
\operatorname{Cov}(X_n,X_{n+k})=R(k).
\]

Эргодичность:

\[
S^{-1}A=A
\Rightarrow
P(A)\in\{0,1\},
\]

или для среднего во втором порядке:

\[
\frac1N\sum_{n=0}^{N-1}X_n
\xrightarrow{L^2}
\mathbb E X_0.
\]

Главная ловушка: стационарность, эргодичность и независимость - разные
свойства.

\newpage

\end{document}
